(a) Write $p(x)=\|x\|$ and $q(x)=\|x\|'$. Lipschitz equivalence means that $mp(x)\le q(x)\le Mp(x)$ for all $x$, with fixed $m,M>0$. These inequalities imply $B_{1/M}\subseteq B'_1\subseteq B_{1/m}$. Conversely, assume the two inclusions with radii $r,R$. For every $s>p(x)/r$, the point $x/s$ lies in $B_r$, so $q(x)<s$ and hence $q(x)\le p(x)/r$. For every $s>q(x)$, the point $x/s$ lies in $B'_1$, giving $p(x)<Rs$ and hence $p(x)\le Rq(x)$. This yields both Lipschitz bounds, including when $x=0$.

(b) Applying the bounds to $x_n-x$ proves that equivalent norms have the same convergent sequences and limits. Conversely suppose the sequence criterion holds. If no upper bound $q\le Mp$ existed, for each positive integer $n$ we could choose $v_n$ with $q(v_n)>n p(v_n)$. Then $u_n=v_n/q(v_n)$ satisfies $p(u_n)<1/n$ while $q(u_n)=1$, contradicting the criterion at zero. Hence $q\le Mp$. Interchanging the norms gives $p\le Cq$ and the conclusion.

(c) The sum $N(x)=p(x)+|\varphi(x)|$ is nonnegative and vanishes only at zero since $p$ is definite. Linearity of $\varphi$ gives absolute homogeneity. The two triangle inequalities add to give the one for $N$. If $N$ were Lipschitz equivalent to $p$, then $|\varphi(x)|\le N(x)\le Cp(x)$ for some $C$. Applying this to differences would give continuity of $\varphi$ everywhere, a contradiction for a discontinuous functional.

(d) Yes, in the usual setting with the algebraic basis theorem. If $V$ is infinite dimensional, successively choose linearly independent unit vectors $e_1,e_2,\ldots$ and extend them to a Hamel basis. Define a linear functional by $\varphi(e_n)=n$ and by zero on the remaining basis vectors. This is well defined since each algebraic basis expansion is finite. The vectors $e_n/n$ tend to zero in the given norm, but $\varphi(e_n/n)=1$, so the functional is discontinuous. Part (c) gives an inequivalent norm, contradicting the hypothesis. For a complex space use a complex-linear functional taking these same values and the complex modulus in the construction.
